% Web専用原稿．独立査読は未了．下記の文献との部分的な照合を実施．
% 参考確認（2026-11-27）：John M. Lee, Introduction to Topological Manifolds, 2nd ed., Springer, 2011, 定理4.68の完備距離空間の場合を確認；Bogachev and Smolyanov, Real and Functional Analysis, 定理1.5.1の閉球の記述を確認．
% 原稿の論証は独自に再構成．書籍全体・全命題との照合は未了．
\section{稠密集合とカテゴリー}
% genron:begin K-000832
% genron:use K-000201
% genron:use K-000001
% genron:use K-000002
% genron:use K-000003
% genron:use K-000004
% genron:use K-000005
% genron:use K-000007
% genron:use K-000017
% genron:use K-000024
% genron:use K-000198
% genron:use K-000199
% genron:use K-000146
% genron:use K-000045
\begin{defn}[疎集合と第一類集合]
位相空間 $X$ の部分集合 $A$ について $\operatorname{Int}\overline A=\varnothing$ となるとき，$A$ は疎であるという．可算個の疎集合の合併として表される集合を第一類集合という．第一類でない集合を第二類集合という．
\end{defn}
% genron:end K-000832
% genron:begin K-000833
% genron:use K-000201
% genron:use K-000001
% genron:use K-000002
% genron:use K-000003
% genron:use K-000004
% genron:use K-000005
% genron:use K-000007
% genron:use K-000017
% genron:use K-000024
% genron:use K-000198
% genron:use K-000199
% genron:use K-000146
% genron:use K-000045
\begin{defn}[Baire空間]
位相空間 $X$ において任意の稠密開集合列 $(U_n)$ に対し $\bigcap_n U_n$ が稠密であるとき，$X$ を Baire 空間という．
\end{defn}
% genron:end K-000833
% genron:begin K-000834
% genron:use K-000146
% genron:use K-000148
% genron:use K-000001
% genron:use K-000002
% genron:use K-000003
% genron:use K-000004
% genron:use K-000005
% genron:use K-000007
% genron:use K-000017
% genron:use K-000024
% genron:use K-000833
% genron:use K-000832
% genron:use K-000198
% genron:use K-000200
% genron:use K-000201
\begin{prop}[Baire性のカテゴリーによる判定]
位相空間が Baire 空間であるための必要十分条件は，空でない開集合が第一類集合に含まれないことである．
\end{prop}
\begin{proof}
疎集合 $A_n$ の閉包の補集合は稠密開集合である．Baire 性を仮定すればこれらの共通部分は稠密なので，空でない開集合と交わる．従ってその開集合は $\bigcup_n A_n$ に含まれない．逆に稠密開集合 $U_n$ の補集合 $X\setminus U_n$ は内部が空な閉集合なので疎である．共通部分が稠密でなければ，これと交わらない空でない開集合があり，その開集合は $\bigcup_n(X\setminus U_n)$ に含まれる．仮定に反するので共通部分は稠密である．
\end{proof}
% genron:end K-000834
\section{完備距離空間のBaire性}
% genron:begin K-000835
% genron:use K-000735
% genron:use K-000100
% genron:use K-000169
% genron:use K-000168
% genron:use K-000001
% genron:use K-000002
% genron:use K-000003
% genron:use K-000004
% genron:use K-000005
% genron:use K-000007
% genron:use K-000017
% genron:use K-000024
% genron:use K-000833
% genron:use K-000769
% genron:use K-000725
% genron:use K-000726
% genron:use K-000727
% genron:use K-000731
% genron:use K-000738
% genron:use K-000060
% genron:use K-000105
% genron:use K-000706
\begin{prop}[Baireのカテゴリー定理]
完備距離空間は Baire 空間である．
\end{prop}
\begin{proof}
完備距離空間 $X$ の稠密開集合列 $(U_n)$ と空でない開集合 $V$ を取る．$V\cap U_1$ は空でない開集合であるから，$x_1$ と $0<r_1<1/2$ を，閉球 $\overline B(x_1,r_1)$ が $V\cap U_1$ に含まれるよう選べる．実際，その開集合に属する点を一つ取り，その点の十分小さい開球が含まれることを利用し，さらに半径を小さくすればよい．

$x_n,r_n$ が選ばれたとき，$B(x_n,r_n)\cap U_{n+1}$ は空でない開集合である．従って
\[
 0<r_{n+1}<2^{-(n+1)},\quad
 \overline B(x_{n+1},r_{n+1})
 \subseteq B(x_n,r_n)\cap U_{n+1}
\]
を満たすものが存在する．選択公理により可能な選択に対する選択函数を固定し，再帰によってこの操作を反復する．$m\geq n$ なら $x_m\in\overline B(x_n,r_n)$ であるから，$p,q\geq n$ なら $d(x_p,x_q)\leq2r_n<2^{1-n}$ である．帰納法から $2^n\geq n+1$ であり，Archimedes性から $2^{1-n}$ は任意に小さくできる．従って $(x_n)$ は Cauchy 列となり，完備性により極限 $x$ を持つ．各閉球は閉集合であり，点列の尾部を含むため $x$ も含む．従って $x\in V\cap\bigcap_n U_n$ である．任意の空でない開集合 $V$ についてこのことが成立するから，共通部分は稠密である．
\end{proof}
% genron:end K-000835
% genron:begin K-000836
% genron:use K-000146
% genron:use K-000148
% genron:use K-000001
% genron:use K-000002
% genron:use K-000003
% genron:use K-000004
% genron:use K-000005
% genron:use K-000007
% genron:use K-000017
% genron:use K-000024
% genron:use K-000835
% genron:use K-000834
% genron:use K-000832
\begin{prop}[閉集合による可算被覆と内部]
空でない完備距離空間 $X$ が閉集合列 $(F_n)$ によって $X=\bigcup_n F_n$ と覆われるならば，ある $F_n$ は空でない内部を持つ．
\end{prop}
\begin{proof}
すべての $F_n$ の内部が空ならば，すべて疎集合である．すると $X$ 自身が空でない開集合でありながら第一類集合となり，Baire性のカテゴリーによる判定に反する．
\end{proof}
% genron:end K-000836
